\usepackage[a4paper,margin=2.2cm,headheight=15pt]{geometry}
\usepackage{amsmath,amssymb,mathtools,amsthm,bm}
\usepackage{enumitem}
\usepackage{xcolor}
\usepackage{fancyhdr}
\usepackage[hidelinks,unicode]{hyperref}
\usepackage[nameinlink,noabbrev]{cleveref}
\allowdisplaybreaks[2]
\definecolor{answerblue}{HTML}{1F4E79}
\definecolor{skipgray}{HTML}{666666}
\setlength{\parindent}{2em}
\setlength{\parskip}{0.25em}
\setlist{itemsep=0.25em,topsep=0.35em}
\pagestyle{fancy}
\fancyhf{}
\fancyhead[L]{泛函分析作业完整解答}
\fancyhead[R]{\nouppercase{\rightmark}}
\fancyfoot[C]{\thepage}
\newtheoremstyle{problemstyle}{0.8em}{0.5em}{\normalfont}{}{\bfseries\color{answerblue}}{.}{0.5em}{}
\theoremstyle{problemstyle}
\newtheorem*{problem}{题干}
\newenvironment{solution}{\par\smallskip\noindent\textbf{解答.}\ }{\hfill$\square$\par\medskip}
\newenvironment{skipped}{\par\smallskip\noindent\color{skipgray}\textbf{本题跳过.}\ }{\par\medskip\color{black}}
\newcommand{\R}{\mathbb R}
\newcommand{\C}{\mathbb C}
\newcommand{\N}{\mathbb N}
\newcommand{\F}{\mathbb F}
\newcommand{\cL}{\mathcal L}
\newcommand{\ol}[1]{\overline{#1}}
\DeclareMathOperator{\Ker}{Ker}
\DeclareMathOperator{\Ran}{Ran}
\DeclareMathOperator{\Span}{span}
\DeclareMathOperator{\dist}{dist}
\DeclareMathOperator*{\esssup}{ess\,sup}
